\section{Discussion}
\label{sec:discussion}

\paragraph{Why does scene-level outperform individual-level?}
The performance gap between DAiSEE ($\kappa \leq 0.10$) and SCB
($\kappa$ up to 0.60) is not simply a function of dataset difficulty.
We hypothesise three contributing factors.  First, \emph{behavioural
salience}: SCB labels derive from overt physical actions (raised hand,
phone use) that are visually unambiguous to a language-grounded model,
whereas DAiSEE labels reflect internal cognitive states that manifest
only subtly in facial microexpressions.  Second, \emph{context richness}:
a classroom scene provides spatial layout, group dynamics, and object
context (\eg, desks, boards) that help ground engagement inference,
while a webcam close-up of a single face offers minimal contextual anchors.
Third, \emph{label construction}: the SCB aggregation scheme maps
explicit physical categories to ordinal scores, creating a more
objective signal than subjective per-clip engagement ratings.

\paragraph{Implications for system design.}
Our findings suggest that deploying VLMs as individual student observers
in their current zero-shot form is premature.  A more productive
paradigm may be \emph{scene-first analysis}: aggregate visual information
at the classroom level, then query a VLM.  This approach aligns better
with how teachers actually monitor classrooms---by scanning the room
rather than studying individual faces.  For privacy reasons, scene-level
analysis is also preferable as it avoids close facial analysis of minors.

\paragraph{Prompt engineering as unreliable compensation.}
The wide prompt sensitivity we observe means that practitioners who tune
prompts on a small validation set risk over-fitting to that distribution.
A prompt that yields 67.9\% accuracy for GPT-4o on SCB could easily
degrade by 28 pp on a different class distribution.  Reliable deployment
requires either task-specific fine-tuning or principled prompt selection
with held-out validation.

\paragraph{Open vs. closed source.}
LLaVA-1.5-7B (4-bit quantised, running locally) matches or outperforms
GPT-4o on DAiSEE, and achieves competitive performance on SCB at zero
API cost and without data leaving the institution.  This is
encouraging for privacy-sensitive educational deployments, though the
comparison is not perfectly controlled given the quantisation difference.

\paragraph{Limitations.}
This study is limited to zero-shot evaluation.  Fine-tuning VLMs on even
modest task-specific data would likely close the gap substantially.
The single-frame assumption discards temporal dynamics known to carry
engagement information.  The SCB engagement labels are derived from
behaviour categories rather than direct human ratings, which may not
fully align with teachers' engagement judgements.  Finally, results on
21 subjects in a webcam setting (DAiSEE) may not generalise to diverse
classroom environments.

\paragraph{Ethical considerations.}
This study relies exclusively on publicly available datasets
(DAiSEE~\cite{gupta2016daisee} and SCB~\cite{wang2023scb}), both
released for academic research under their respective
terms of use.  No new human data was collected, and no individuals
are re-identified or tracked.  DAiSEE was originally collected under
informed consent with institutional ethical oversight; our use is
restricted to zero-shot inference on pre-existing anonymised clips.
The involvement of minors in classroom imagery underscores the
importance of privacy-preserving design: our findings---particularly
the GPT-4o safety refusals and the superior performance of
scene-level over individual-level analysis---reinforce that
aggregated, non-facial approaches are both more tractable and more
ethically appropriate for real-world educational deployment.
